\subsection{Temperature}

In the final step of GPT decoding, a logit vector is normalized into a probability distribution over the vocabulary using the Softmax function:
\begin{equation}
	P(z_i) = \frac{e^{z_i}}{\sum_j{e^{z_j}}}.
\end{equation}

Each value in the logit vector \((z_1, z_2, \dots, z_V)\) represents ``how likely the corresponding word is to be selected as the next token,'' where \(V\) denotes the vocabulary size. The Softmax function converts these logits to a probability distribution, assigning higher probabilities to words with larger logits. GPT then samples the next token from this distribution.

Sometimes, we want the model be more creative and generate words or sentences that are less common. Mathematically, this means making the final distribution flatter, so that lower-probability words are more likely to be sampled. To achieve this, we introduce a temperature parameter \(T\), which divides all logits before they are passed into the Softmax function:
\begin{equation}
	P(z_i) = \frac{e^{z_i/T}}{\sum_j{e^{z_j/T}}}.
\end{equation}

When \(T > 1\), the differences between logits become smaller, which makes the resulting probability distribution flatter. As a result, the most likely token becomes less dominant, and less likely tokens receive higher probabilities. In contrast, when \(0 < T < 1\), the differences between logits are amplified, making the distribution sharper. In this case, the most likely token becomes even more dominant.

For example, suppose we have a vocabulary of size 5: \texttt{(cat, dog, rabbit, horse, frog)}, with logits \((2.0, 9.5, 4.0, 1.5, 0.5)\). The resulting probability distributions under different temperatures are shown in Table~\ref{tab:probability_distribution_over_the_vocabulary}.

\begin{table}[h]
	\centering
	\begin{tabular}{lcccc}
		\toprule
		Word   & Logit & \(P_{T = 1}\) & \(P_{T = 2}\) & \(P_{T = 0.8}\) \\
		\midrule
		cat    & 2.0   & 0.0005        & 0.0211        & 0.0001          \\
		dog    & 9.5   & 0.9949        & 0.8954        & 0.9988          \\
		rabbit & 4.0   & 0.0041        & 0.0572        & 0.0010          \\
		horse  & 1.5   & 0.0003        & 0.0164        & 0.0001          \\
		frog   & 0.5   & 0.0001        & 0.0099        & 0.0000          \\
		\bottomrule
	\end{tabular}
	\caption{Probability distribution over the vocabulary for different temperatures.}
	\label{tab:probability_distribution_over_the_vocabulary}
\end{table}
